\documentclass[11pt]{article}

\usepackage[T1]{fontenc}
\usepackage{lmodern}
\usepackage[margin=1in]{geometry}
\usepackage{microtype}
\usepackage{amsmath,amssymb,amsthm,mathtools}
\usepackage{bm}
\usepackage{graphicx}
\graphicspath{{figures/}}
\usepackage{booktabs}
\usepackage{subcaption}
\usepackage{xcolor}
\usepackage[numbers,sort&compress]{natbib}
\usepackage[colorlinks,allcolors=blue!60!black]{hyperref}
\usepackage[capitalize,noabbrev]{cleveref}
\usepackage{thmtools}
% Theorem environments, numbered together (Theorem 1, Lemma 2, ...).
% \declaretheorem (thmtools) keeps cleveref's names right for shared counters:
% \cref{def:x} -> "Definition 1", not "Theorem 1".
\declaretheorem[style=plain]{theorem}
\declaretheorem[style=plain, sibling=theorem]{lemma}
\declaretheorem[style=plain, sibling=theorem]{proposition}
\declaretheorem[style=definition, sibling=theorem]{definition}
\declaretheorem[style=definition]{assumption}
\declaretheorem[style=remark, sibling=theorem]{remark}

% Sets and operators.
\newcommand{\R}{\mathbb{R}}
\newcommand{\N}{\mathbb{N}}
\newcommand{\E}{\mathbb{E}}
\DeclareMathOperator*{\argmin}{arg\,min}
\DeclareMathOperator*{\argmax}{arg\,max}
\DeclarePairedDelimiter{\abs}{\lvert}{\rvert}
\DeclarePairedDelimiter{\norm}{\lVert}{\rVert}

% Vectors and matrices: bold, e.g. \vect{x}, \mat{A}.
\newcommand{\vect}[1]{\bm{#1}}
\newcommand{\mat}[1]{\bm{#1}}

% Inline note for unfinished parts; remove all before sharing.
\newcommand{\todo}[1]{\textcolor{red}{[TODO: #1]}}


\title{TITLE}
\author{AUTHOR}
\date{\today}

\begin{document}
\maketitle

\begin{abstract}
	% One paragraph: the problem, why it matters, and where things stand.
\end{abstract}

\section{Problem}
\label{sec:problem}
% What exactly is unsolved, stated precisely enough that someone could tell
% whether a proposed answer solves it. Inputs, outputs, success criteria.

\section{Notation}
\label{sec:notation}
% Every symbol used later, defined once.

\section{Background}
\label{sec:background}
% Prior work and why it does not already solve the problem. Cite with
% \citet{key} or \citep{key}; entries go in refs.bib.

\section{Approach}
\label{sec:approach}
% The proposed idea, its key equations, and why it should work.

\section{Open questions}
\label{sec:open}
% What is unknown or risky, and how each question could be answered.

\bibliographystyle{plainnat}
\bibliography{refs}

\end{document}
